\section*{文字と説明について}
特に断らなければ，変数は実数を表す。
$\log x$は自然対数，$\e$はその底とし，三角関数の角は弧度法で測る。
$\mathbb{N}$を用いる場合は正の整数の集合とする。
不定積分に現れる$C$は積分定数であり，式ごとに値が変わってよい。
分母が零になる点などを除いて区間が分かれる場合には，それぞれの区間で定数を選ぶ。

定積分の$\dd x$は，どの変数について積分しているかを示す。
置換の途中で$\dd u=g'(x)\dd x$と書くときには，置換積分の公式を簡潔に表している。
この記号だけを根拠に，どのような操作でも分数のように行えると考えないようにしたい。

重要な公式は，その公式を使う理由と導出を本文に示した。
一方，連続関数の積分可能性など，実数の性質にまで戻る証明は，採用する事実として明示している。
発想を得るための近似と，結果が正しいことを確かめる議論も分けてある。
例題と章末問題は本書の説明のために設定したものであり，特定の入試問題の転載ではない。
章末問題と総合演習には，巻末にすべて解答を付けた。
